% ============================================================
\section{Formal Invariants and Proofs}
\label{sec:formal-invariants}

This section states and proves floor monotonicity by case analysis over the protocol's operation set, establishes the integer-arithmetic redemption lemma, states the redemption-fairness corollary, a closure result under arbitrary interleaving, and a directional conjecture inverting the LVR framing of \citet{MilionisMoallemiRoughgarden2023}.

\paragraph{State and operations.} System state:
\begin{equation}
\sigma = (\Tres, \Sup, \Lt, \Ls, \Phit, \Phis),
\label{eq:state}
\end{equation}
with $\Tres$ the treasury balance, $\Sup$ total token supply, $(\Lt, \Ls)$ LP reserves, and $(\Phit, \Phis)$ the raw unrealized fee mass on the LP position awaiting the next \texttt{collectFees}.\footnote{The Uniswap v4 implementation does not store $\Phit, \Phis$ as raw token amounts; instead, it maintains per-currency Q128.128 fee-growth accumulators \texttt{feeGrowthGlobalX128} into which each swap's input-side fee is folded, with the position's claim computed as $K_{\text{real}} = (\text{accumulator delta}) \cdot \text{liquidity} / 2^{128}$ at realization. Throughout this paper we use the raw-token framing $(\Phit, \Phis)$ for dimensional clarity of the proofs; the v4 Q128.128 realization is an implementation detail that the algebra abstracts away.} A token-input swap of $N_{\text{in}}$ adds $f_s N_{\text{in}}$ to $\Phit$; a stable-input swap of $X_{\text{in}}$ adds $f_b X_{\text{in}}$ to $\Phis$. \texttt{collectFees} sets $(\Kreal, \Ureal) = (\Phit, \Phis)$ and resets both to zero. Derived quantities: $\Floor = \Tres/\Sup$ (integer-faithful per~\eqref{eq:floor-precision}) and $\Spot = \Ls/\Lt$. The constant-product $\kInv = \Lt \cdot \Ls$ is preserved by every swap modulo the asymmetric fee; exact because the position is full-range.

The protocol's \emph{operation set} $\mathrm{Ops}$ is the disjoint union of seven operation classes:
\begin{align}
\mathrm{Ops} = \{\;
  &\text{RevToTreasury}(X),
  \;\text{BuyAndBurn}(X),
  \;\text{Redeem}(N), \notag \\
  &\text{LPBuy}(X_{\text{in}}),
  \;\text{LPSell}(N_{\text{in}}),
  \;\text{Transfer}(N, a \to b),
  \;\text{CollectFees}() \;\}.
\label{eq:ops}
\end{align}
The first two are the halves of \texttt{routeRevenue}; the next three are holder-facing; the sixth is ERC-20 transfer; the seventh realizes fees. No other entry point exists in the bytecode.

\begin{lemma}[Integer-Arithmetic Redemption]
\label{lem:integer-redemption}
Let $(\Tres, \Sup) \in \mathbb{Z}_{> 0}^2$ with $N \in \mathbb{Z}_{> 0}$ and $N \leq \Sup$. Let the on-chain redemption payout be the integer
\begin{equation}
P = \left\lfloor \frac{N \cdot \Tres}{\Sup} \right\rfloor, \qquad
\text{and write } r = (N \cdot \Tres) \bmod \Sup \in \{0, 1, \ldots, \Sup-1\}.
\label{eq:redeem-payout}
\end{equation}
Let $\Tres' = \Tres - P$ and $\Sup' = \Sup - N$. Then either $\Sup' = 0$---in which case $N = \Sup$ forces $r = 0$ and $P = \Tres$, so $\Tres' = 0$ exactly---or $\Sup' > 0$ and
\begin{equation}
\frac{\Tres'}{\Sup'} \;=\; \frac{\Tres}{\Sup} + \frac{r}{\Sup \cdot \Sup'} \;\geq\; \frac{\Tres}{\Sup}.
\label{eq:integer-redemption}
\end{equation}
In particular, $\Floor_{\mathrm{new}} \geq \Floor_{\mathrm{old}}$, with equality if and only if $\Sup \mid N \cdot \Tres$.
\end{lemma}

\begin{proof}
By definition of integer division, $N \cdot \Tres = \Sup \cdot P + r$ with $0 \leq r < \Sup$. Then
\[
\Tres - P \;=\; \Tres - \frac{N \cdot \Tres - r}{\Sup} \;=\; \frac{\Tres(\Sup - N) + r}{\Sup} \;=\; \frac{\Tres \cdot \Sup' + r}{\Sup}.
\]
If $\Sup' > 0$, dividing by $\Sup'$ gives $\Tres'/\Sup' = \Tres/\Sup + r/(\Sup \cdot \Sup')$, which is~\eqref{eq:integer-redemption}. The right-hand correction term is non-negative; it is zero exactly when $r = 0$, i.e.\ when $\Sup$ divides $N \cdot \Tres$ exactly. If $\Sup' = 0$, then $N = \Sup$ and $N \cdot \Tres = \Sup \cdot \Tres$ is divisible by $\Sup$, so $r = 0$ and $P = \Tres$; hence $\Tres' = \Tres - P = 0$ exactly. This is the terminal redemption that drains the contract; the floor at this point is undefined, and the contract reverts on the next \texttt{floorPrice}() call via the \texttt{SupplyExhausted} error.
\end{proof}

Lemma~\ref{lem:integer-redemption} sharpens ``invariant up to truncation'' to an exact inequality: integer truncation strictly favors the treasury.

\begin{theorem}[Floor Monotonicity]
\label{thm:floor-monotonicity}
Let $\sigma$ be any reachable state of the protocol and let $\text{op} \in \mathrm{Ops}$ be any operation applicable in $\sigma$. Let $\sigma' = \text{op}(\sigma)$. Then
\begin{equation}
\Floor(\sigma') \geq \Floor(\sigma),
\label{eq:floor-monotone}
\end{equation}
with the equality and strict-inequality cases as enumerated in the proof.
\end{theorem}

\begin{proof}
We argue by cases over the seven operation classes of~\eqref{eq:ops}. In each case we write $(\Tres, \Sup) \to (\Tres', \Sup')$ for the treasury and supply transition and compute the resulting floor.

\textbf{Case 1 --- $\text{RevToTreasury}(X)$ with $X \geq 0$.} The router moves $X$ stable to the treasury. State transition: $(\Tres + X, \Sup)$. Floor: $(\Tres + X)/\Sup \geq \Tres/\Sup$, with strict inequality iff $X > 0$.

\textbf{Case 2 --- $\text{BuyAndBurn}(X)$ with $X \geq 0$.} The router uses $X$ stable to execute a buy on the LP, receiving and burning $Y(X) = \Lt - \kInv/(\Ls + (1 - f_b) X)$ tokens, where $f_b = 0.001$ is the buy-side hook fee. State transition: $(\Tres, \Sup - Y(X))$. Floor: $\Tres/(\Sup - Y(X)) \geq \Tres/\Sup$, with strict inequality iff $Y(X) > 0$, which holds for any $X > 0$ provided $\Ls$ is finite and positive.

\textbf{Case 3 --- $\text{Redeem}(N)$ with $1 \leq N \leq \Sup$.} By Lemma~\ref{lem:integer-redemption}, $\Floor_{\mathrm{new}} \geq \Floor_{\mathrm{old}}$, with equality iff $\Sup \mid N \cdot \Tres$. For the abstract rational ideal $\Floor = \Tres/\Sup$, the floor is preserved exactly; under integer truncation, the floor strictly rises by a sub-unit residual that accrues to the treasury.

\textbf{Case 4 --- $\text{LPBuy}(X_{\text{in}})$ from an external user.} The user sends $X_{\text{in}}$ stable to the LP and receives tokens out. The LP reserves change; the treasury and total supply are untouched. State transition: $(\Tres, \Sup)$, so $\Floor' = \Floor$. Equality holds in every instance.

\textbf{Case 5 --- $\text{LPSell}(N_{\text{in}})$ from an external user, with the hook's 3\% fee on input.} Let $f_s = 0.03$ be the sell-side fee. The hook deducts $f_s \cdot N_{\text{in}}$ from the input before the curve, folding the deducted amount into $\Phit$. The remaining $(1 - f_s) N_{\text{in}}$ swaps through the curve. State transition: $(\Tres, \Sup, \Lt + (1 - f_s) N_{\text{in}}, \Ls - \Delta\Ls, \Phit + f_s N_{\text{in}}, \Phis)$, where $\Delta\Ls = \Ls - \kInv/(\Lt + (1 - f_s) N_{\text{in}})$ is the stable amount sent to the seller. The treasury and total supply are unchanged, so $\Floor' = \Floor$ at the moment of the swap. The fee tokens reside in the v4 fee accumulator and are not yet on the protocol's books; they are realized in Case~7.

\textbf{Case 6 --- $\text{Transfer}(N, a \to b)$.} An ERC-20 balance transfer between addresses $a$ and $b$. Treasury, supply, and LP reserves are untouched. $\Floor' = \Floor$ trivially.

\textbf{Case 7 --- $\text{CollectFees}()$.} The function realizes the accrued unrealized fee mass as $(\Kreal, \Ureal) = (\Phit, \Phis)$ and resets both accumulators to zero (cf.\ the Q128.128 accumulator semantics noted earlier in this section: the v4 implementation extracts the same quantities scaled by the position's liquidity). The token-side $\Kreal$ is split into $K_b = \lfloor 0.0025 \Kreal \rfloor$ (caller bounty, remains in circulation) and $K_{\text{burn}} = \Kreal - K_b$ (burned). The stable-side $\Ureal$ is split into $U_b = \lfloor 0.0025 \Ureal \rfloor$ (caller bounty) and $U_{\text{treas}} = \Ureal - U_b$ (added to treasury). State transition: $(\Tres + U_{\text{treas}}, \Sup - K_{\text{burn}}, \Lt, \Ls, 0, 0)$. New floor:
\[
\Floor' = \frac{\Tres + U_{\text{treas}}}{\Sup - K_{\text{burn}}} \;\geq\; \frac{\Tres}{\Sup} = \Floor,
\]
with strict inequality iff $U_{\text{treas}} > 0$ or $K_{\text{burn}} > 0$, which holds whenever any swap has occurred since the previous \texttt{collectFees} call.

\textbf{Exhaustiveness.} The seven cases exhaust $\mathrm{Ops}$ as defined in~\eqref{eq:ops}. The protocol bytecode exposes no other externally callable function that modifies $(\Tres, \Sup, \Lt, \Ls)$: the token has \{\texttt{redeem}, \texttt{burn} (restricted), \texttt{transfer}, \texttt{transferFrom}, \texttt{approve}, \texttt{permit}\}; the treasury has \{\texttt{payRedemption} (restricted)\}; the router has \{\texttt{routeRevenue} (restricted)\}; the hook has \{\texttt{beforeSwap}, \texttt{collectFees}\}. Approval and permit do not modify the state tuple $\sigma$; \texttt{beforeSwap} runs as part of a swap (subsumed by Cases~4 and~5); the restricted functions can be invoked only by the listed callers and so are subsumed by the protocol-initiated operations enumerated above. No other path can reach $\sigma$.\footnote{The hook's \texttt{initializePool} is a one-shot genesis call that transitions $(\Lt, \Ls)$ from zero to the genesis seed (Section~\ref{sec:protocol}) and reverts on every subsequent call; \texttt{unlockCallback} is \texttt{onlyPoolManager}-gated and its $\sigma$-mutations are subsumed by Cases~2, 4, 5, and~7 (router buy-and-burn and AMM swaps re-enter the hook through this callback). Neither path introduces a state transition outside $\mathrm{Ops}$.}

Combining all cases: $\Floor(\sigma') \geq \Floor(\sigma)$ holds for every $\text{op} \in \mathrm{Ops}$, with strict inequality in Cases~1, 2, 7 (and weakly in Case~3 under the abstract ideal, strictly under integer arithmetic with $r > 0$), and equality in Cases~4--6. $\qed$
\end{proof}

\begin{corollary}[Closure under Arbitrary Interleaving]
\label{cor:closure}
For any finite sequence $\text{op}_1, \text{op}_2, \ldots, \text{op}_m \in \mathrm{Ops}$ applied to a reachable state $\sigma$, the resulting state $\sigma_m = \text{op}_m \circ \cdots \circ \text{op}_1(\sigma)$ satisfies $\Floor(\sigma_m) \geq \Floor(\sigma)$.
\end{corollary}

\begin{proof}
By induction on $m$. The base case $m = 0$ is trivial. For the inductive step, $\Floor(\sigma_m) \geq \Floor(\sigma_{m-1}) \geq \Floor(\sigma)$ by Theorem~\ref{thm:floor-monotonicity} and the induction hypothesis. $\qed$
\end{proof}

Corollary~\ref{cor:closure} absorbs intra-block ordering ambiguity: no adversarial reordering or insertion can drop the floor, so the invariant is robust to all single-block reorderings.

\begin{theorem}[Redemption Fairness, No First-Mover Advantage]
\label{thm:redemption-fairness}
Let two holders, $a$ and $b$, simultaneously submit redemption transactions for $N_a, N_b$ tokens against a state $\sigma$. Let any miner-/builder-chosen ordering execute these (and possibly other in-protocol operations) sequentially, yielding payouts $P_a^{(\pi)}, P_b^{(\pi)}$ under ordering $\pi$. Then for every ordering $\pi$ in which $a$ executes before $b$:
\begin{equation}
P_a^{(\pi)} / N_a \;\leq\; P_b^{(\pi)} / N_b.
\label{eq:fairness}
\end{equation}
That is, the per-token payout for the later redeemer is weakly higher than that of the earlier redeemer; no first-mover advantage exists in floor terms.
\end{theorem}

\begin{proof}
The per-token payout for a redemption is the floor evaluated at that redemption's block-state. By Theorem~\ref{thm:floor-monotonicity} Case~3 (and Lemma~\ref{lem:integer-redemption} for the integer-arithmetic refinement), $a$'s redemption preserves or raises the floor; by Corollary~\ref{cor:closure}, any intervening operations preserve or raise it further. Therefore $b$'s redemption evaluates the floor at a value at least as high as $a$'s did, establishing~\eqref{eq:fairness}. $\qed$
\end{proof}

Theorem~\ref{thm:redemption-fairness} concerns the floor, not absolute amounts; a holder who could sell at spot above the floor accepts an opportunity cost by redeeming.

\paragraph{Conjecture and LVR.} The LVR framework of \citet{MilionisMoallemiRoughgarden2023} characterizes the cost an AMM LP incurs against informed arbitrageurs. For the protocol-owned LP in M\textsuperscript{2}, this framing appears to invert. We state it as a conjecture; a quantitative identity requires a unit reconciliation between LVR's stochastic-calculus formulation and M\textsuperscript{2}'s discrete-event floor framing.

\begin{conjecture}[LVR / Floor-Capture Identity]
\label{conj:lvr-capture}
Let $\mathrm{LVR}_{[t_0, t_1]}$ denote the cumulative loss-versus-rebalancing of the protocol-owned LP position over an interval $[t_0, t_1]$, in the sense of \citet{MilionisMoallemiRoughgarden2023}. Under the asymmetric-fee mechanism of Section~\ref{sec:protocol} with $f_b = 0.001$ on buys and $f_s = 0.03$ on sells, the LVR-equivalent value of the protocol-owned LP position is recovered as floor growth via the Case~5 plus Case~7 composition of Theorem~\ref{thm:floor-monotonicity}, modulo a leakage of at most $0.25\%$ per side at \texttt{collectFees} realization. In particular, on any interval $[t_0, t_1]$ over which \texttt{collectFees} is invoked at least once,
\begin{equation}
\Delta \Floor_{[t_0, t_1]} \cdot \Sup(t_1) \;\gtrsim\; (1 - 0.0025) \cdot \mathrm{LVR}^{\text{fee-attributable}}_{[t_0, t_1]},
\label{eq:lvr-capture}
\end{equation}
where $\Delta \Floor_{[t_0, t_1]}$ is the cumulative floor increase attributable to swap-induced fee growth realized via \texttt{collectFees} in the same interval, and the right-hand side is the portion of cumulative LVR that the asymmetric fee mechanism converts into protocol-realizable fee growth.
\end{conjecture}

\paragraph{Remark (status).} Conjecture~\ref{conj:lvr-capture} is directional; promotion to a theorem is an open problem (P2). The directional content suffices for Section~\ref{sec:related-work}'s structural inversion (iv): the protocol-owned LP converts LVR from a loss-to-LP into floor growth captured by the holder base.
